% Student notes for Herlock from the second outside review (2026-09-22), all checked against our
% code/results. Each block goes right AFTER the anchor line quoted above it. No blank lines inside \vg{}.

% --- N1: after  \end{keywords}
\vg{\textbf{[Students] Submission logistics.} (1) Deadline: please confirm on the CMS submission page; a second reviewer reads it as 07:00 Thu 24 Sep, New York time, earlier than the 07:59 we had. (2) The SPS policies require a \emph{Compliance with Ethical Standards} statement ``irrespective of whether ethical approval was needed''; it may go on page 5. We use IMDb-Wiki faces and Adult census records; their example wording: ``This research study was conducted retrospectively using human subject data made available in open access by [source]. Ethical approval was not required as confirmed by the license attached with the open access data.'' (3) Gietl--Reffel is \emph{Metrika}, vol.~76, no.~6, pp.~783--798, 2013, not \emph{Statistics} 47(6). (4) We re-uploaded figs/rare\_group\_fit.pdf: the Adult x-axis now reads ``sampling weights $\propto p^\beta$'' instead of ``observation rate $r\propto p^\beta$''.}

% --- N2: after  $0\le f_i\le M$ (cf.\ \cite{rahimi2025agnostic}). \hfill$\square$
\vg{Unnormalized weights escape $\Reach$: Horvitz--Thompson weights $p_i/r_i$ on the present groups, not normalized, give $\E\hat g^t=\sum_ir_i(p_i/r_i)s_i=\nabla F_p$ whenever all $r_i>0$, feasible or not. So ``every rule'' (abstract) and ``no rule is unbiased'' (Sec.~3) hold for normalized rules only. On the two infeasible headline cells availability is uniform, so the fixed multiplier $(N/K)p_i$ \emph{is} Horvitz--Thompson: it gives $0.670$ and $8504$ in Table~\ref{tab:overall}. Its second moment reaches $G^2\max_j\sum_{i\in A_j}p_i/r_i=15.6\,G^2$ on Adult $\beta=0$, while Theorem~\ref{thm:rate} keeps $G^2$.}

% --- N3: after  \end{assumption}
\vg{Proposition~\ref{prop:surrogate} and Theorem~\ref{thm:rate} also need $A^t$ drawn i.i.d.\ from $q$ at each step, independent of $\theta^{t-1}$ (the code does this). Epoch-wise sampling without replacement, common in streaming or sharded pipelines, violates it.}

% --- N4: after  that reaches $p$ when this is possible, and the best one when it is not.
\vg{Eq.~\eqref{eq:mono} is an aggregate statement; for $N>2$ an under-weighted group's risk can fall. Example: $f_1=(\theta-1)^2$, $f_2=(\theta+1)^2$, $f_3=(\theta-0.9)^2$; $p$ uniform gives $\theta^\star=0.30$, $w=(0.6,0.3,0.1)$ gives $\theta^\star=0.39$, and group 3 is under-weighted yet $f_3$ drops from $0.36$ to $0.26$. The two-group argument in Sec.~1 is fine; the same caveat applies to ``their risks rise'' in Sec.~4.}

% --- N5: after  weightings decides the sign; \eqref{eq:gain} does.
\vg{(i) $\hat p$ is not always truncated \emph{at} $r_i$: in our Adult $\beta=0$ plan ($\|p-\hat p\|_1=0.912$) the Asian groups get $\hat p=0.0282$ and the Black groups $0.0245$, below $r\approx0.030$, because the Hall constraint on the 15 non-White groups binds (their $\hat p$ sums to $0.3895=1-\binom{85}{3}/\binom{100}{3}$). So ``exactly the groups with $p_i>r_i$'' in Sec.~4 is ``at least''. (ii) $\Phi(\hat p)$ and $\Phi(\tilde p)$ need the weighted optima over all groups, i.e.\ full data access, so $\Delta_{\mathrm{floor}}$ is fixed by the instance but not computable from $(p,q)$ alone ($\nu$ and the feasibility check are). The residual gap is not universal: on Adult $\beta=0$ GAVOT's residual is smaller than blind's ($0.0100$ vs $0.0200$). (iii) The $DG/\sqrt{T}$ bound is on $F_w$, not on $\varepsilon_T(w)$, which has no rate and can be negative.}

% --- N6: after  grows as $\alpha$ and $\gamma$ shrink.
\vg{On Adult $\beta=0$ the premise of \eqref{eq:certificate} needs $\gamma\le\min_ip_i/\hat p_i\approx0.33$ and $\alpha\le(1-\gamma)/(\rho-\gamma)$ with $\rho\approx6.7$, i.e.\ $\alpha\lesssim0.11$--$0.15$, so the reported $(0.2,0.7)$ in Sec.~5 is outside the certified region. Risk-neutral transport already has $F_p\le F_{\hat p}+\tfrac M2\|p-\hat p\|_1$.}

% --- N7: after  of order $p_i/r_i$; it is a guarantee, not a way to lower the average.
\vg{Smaller items, all checked against the runs. (a) Sign flip: the gains at $\beta=1.5$ ($t=1.1$) and $\beta=2$ ($t=-0.1$) are both ties; the residual ranges $1.4$--$2.3$. (b) Table~\ref{tab:overall}: two Adult skews share $\nu=0.60$ ($\beta=0,\frac14$) and two share $\nu=0$ ($\beta=\frac34,1$); the rows shown are $\beta=0$ and $\beta=1$. (c) Aligned Adult: GAVOT also beats blind ($t=7.1$); there $\|p-\tilde p\|_1=0.108$, and the closeness comes from $\Delta_{\mathrm{floor}}=2.6\cdot10^{-4}$. (d) $t=5.8$ and $2.6$ are Adult $\nu=0.60$ and IMDb $\nu=0.31$. (e) Adult wins at every skew even at constant stepsize ($t\ge6.3$). (f) At $\beta=0$ GAVOT ends $2.6$ above its floor, and in Fig.~\ref{fig:severity} the win stops while the floors still differ by $2.6$. (g) Fig.~\ref{fig:rare} shows the whole sweep, not only Table~\ref{tab:overall}'s runs; ``prevalence availability'' is not defined; at $\beta=1.5$ $F_p$ gains while the top identities lose. (h) $F_p$ is the training loss. (i) Notation: $g^t_k$ (label) vs $g^t_i$ (subgradient), $G$ (group) vs $G$ (gradient bound), $\hat{\mathcal Y}$ undefined; $J$ can be exponential in $N$.}
